\section{Conclusion}
\label{sec:conclusion}

Tabular agentic question answering is a control-plane problem. Each candidate source represents a multi-turn execution commitment, and the failure modes that dominate single-agent runs --- schema thrash, zero-row loops, file-handle errors, source drift, premature submit, and context pollution --- arise from the absence of explicit commitment, scoping, and stopping primitives in the runtime. Tabular MOMO is a contract-driven runtime control plane for this problem: a planner orchestrates without touching raw data tools; bounded search and inspect subagents execute explicit contracts; runtime guards enforce file-reference and source-family discipline; and a compact-return protocol with enumerated \texttt{missing\_outputs} surfaces partial progress structurally rather than through raw traces. MOMO is implemented as a modular extension over a baseline \texttt{DataLakeAgent} via five generic extension hooks, exposing \texttt{results}, \texttt{cot}, \texttt{sprint}, and \texttt{delegation} as separable control-plane levers. The harness evaluates at the workflow level --- tool waste, context pollution, source-switch churn, partial-success handling, and cost --- and maps each diagnosed failure mode to the addressing control-plane mechanism. The contribution is not that multi-agent systems work; it is that an execution-grounded control plane with explicit contracts and runtime guards exposes and reduces a class of workflow-level failures that text-centric Agentic RAG cannot.
